

\part{Anhang}
\chapter{Glossar --- Begriffserklärungen}
{\mdseries
\label{bkm:LexComputertomographie}Computertomographie. Abk.: CT. }

{\mdseries
CT. ${\rightarrow}$Computertomographie.}

{\mdseries
Kernspintomographie. ${\rightarrow}$Magnetresonanztomographie}

{\mdseries
Magnetresonanztomographie. Syn.: Kernspintomographie. Abk.: MRT}

{\mdseries
MRT. ${\rightarrow}$Magnetresonanztomographie}

{\mdseries
NAW. ${\rightarrow}$Notarztwagen.}

{\mdseries
Notarztwagen. Abk.: NAW.}

{\mdseries
Röntgen. }

\chapter{Impressum, Hinweise}
\begin{multicols}{2}
{\bfseries Arbeitsgemeinschaft Arbeits- und Ausbildungsstandards für den
Sanitätsdienst} -- \AASS

\emph{(ARGE AASS)}

Anschrift: 

Engerthstraße 146\,/\,8\,/\,13 \\ A-1200 Wien -- AUSTRIA

E-Mail: office{\includegraphics[height=1em]{./images/shared/cmr-at.png}}aass.at
%{\Huge @}

Homepage: \url{http://www.aass.at}

\begin{description}
\item[] 



\end{description}

\paragraph*{Hinweise}~

Wie jede Wissenschaft ist die Medizin ständigen Ent\-wicklungen
unterworfen. Soweit fachliche Angaben ge\-macht werden, darf der Leser
zwar darauf ver\-trau\-en, dass die Autoren gro{\ss}e Sorg\-falt
da\-rauf verwandt haben, dass diese An\-ga\-ben dem Wissens\-stand bei
Ferti\-gstel\-lung des Werkes entspricht. Für Angaben, spe\-ziell
über Dosierungsanweisungen und Ap\-plikations\-formen, kann jedoch
keine Ge\-währ übernommen werden. Jeder Leser ist angehalten, durch
sorg\-fältige Prü\-fung der Literatur, der Lehrmeinung so\-wie der
Beipackzettel der ver\-wendeten Präparate und gegebenenfalls nach
Konsultation eines Spezialisten fest\-zustellen, ob die gegebene
Aussage oder Empfehlung für Do\-sierungen oder die Beachtung von
Kon\-tra\-indi\-ka\-tionen gegenüber der Angabe in dieser
Mit\-tei\-lung abweicht. Jede Dosierung oder Ap\-pli\-kation erfolgt
auf eigene Gefahr des Be\-nutzers. Die ARGE AASS appelliert an jeden Le\-ser
ihm etwa auffallende Un\-ge\-nau\-igkei\-ten umgehend
mitzuteilen.

Geschützte Waren\-namen und Wa\-ren\-zei\-chen werden nicht besonders
kenntlich ge\-macht. Aus dem Fehlen eines solchen Hin\-weises kann also
nicht ge\-schlossen werden, dass es sich um einen freien Waren\-namen
handelt.
\end{multicols}


